\documentclass[12pt,a4paper]{article}
\usepackage[a4paper,margin=20mm]{geometry}
\usepackage[T2A]{fontenc}
\usepackage[utf8]{inputenc}
\usepackage[russian]{babel}
\usepackage{amsmath,amssymb,mathtools}
\usepackage{array,booktabs,longtable}
\usepackage{xcolor}
\usepackage{hyperref}
\usepackage{tikz}
\hypersetup{colorlinks=true,linkcolor=blue!60!black,urlcolor=blue!60!black}
\setlength{\parindent}{0pt}
\setlength{\parskip}{0.6em}
\renewcommand{\arraystretch}{1.25}

\begin{document}

\begin{titlepage}
\thispagestyle{empty}
\begin{center}
\vspace*{0.28\textheight}
{\LARGE\bfseries Ревью домашней работы\par}
\vspace{1.5cm}
{\large Автор: Лёвин Артём Александрович\par}
\vspace{0.7cm}
{\large 13 сентября 2026 года\par}
\vfill
\begin{tikzpicture}[x=1cm,y=1cm]
\draw[line width=0.8pt] (0,0) .. controls (2.2,0.7) and (4.5,-0.7) .. (6.7,0);
\end{tikzpicture}
\end{center}
\end{titlepage}

\newpage
\section*{Сводная таблица}
\small
\setlength{\tabcolsep}{3pt}
\begin{longtable}{@{}>{\centering\arraybackslash}p{0.06\linewidth} p{0.22\linewidth} p{0.23\linewidth} p{0.18\linewidth} p{0.16\linewidth}@{}}
\toprule
\textbf{№} & \textbf{Тема} & \textbf{Суть ошибки} & \textbf{Ваш ответ} & \textbf{Правильный ответ}\\
\midrule
\endfirsthead
\toprule
\textbf{№} & \textbf{Тема} & \textbf{Суть ошибки} & \textbf{Ваш ответ} & \textbf{Правильный ответ}\\
\midrule
\endhead
\hyperref[task:6]{6} & Единственность плоскости по двум пересекающимся прямым & Из общего пункта двух плоскостей сделан вывод об их совпадении; нужная теорема о единственности плоскости по двум пересекающимся прямым не применена. & Сделан вывод, что прямые лежат в одной плоскости. & Плоскости $\alpha$ и $\beta$ совпадают.\\
\addlinespace
\hyperref[task:9]{9} & Единственность плоскости по прямой и внешней точке & В доказательстве сказано, что две точки на прямой задают единственную плоскость. Двух точек для этого недостаточно; нужно использовать точку $M\notin a$. & $\alpha=\beta$, но обоснование неполно. & $\alpha=\beta$.\\
\addlinespace
\hyperref[task:10]{10} & Скрещивающиеся прямые и плоскости & Записаны данные и рисунок, но доказательство различия плоскостей не приведено. & Доказательство отсутствует. & $\alpha\ne\beta$.\\
\bottomrule
\end{longtable}
\normalsize

\newpage
\thispagestyle{empty}
\vspace*{0.42\textheight}
\begin{center}
{\LARGE\bfseries Разбор ошибок}
\end{center}
\newpage

\phantomsection\label{task:6}
\section*{Задание 6}

\textbf{Текст задания.}
Две плоскости $\alpha$ и $\beta$ содержат пересекающиеся прямые $p$ и $q$. Докажите совпадение плоскостей.

\textbf{Ваше решение.}
В записи использовано рассуждение примерно следующего вида: две пересекающиеся прямые задают одну точку; точка $p\cap q$ принадлежит и $\alpha$, и $\beta$, после чего сделан вывод, что прямые лежат в одной плоскости.

\textbf{Ваш ответ.}
Явного вывода $\alpha=\beta$ в доказательстве нет; записано, что прямые лежат в одной плоскости.

\textcolor{red}{\textbf{Ошибка:}} общего пункта двух плоскостей недостаточно, чтобы доказать их совпадение.

\textbf{Где именно возникает ошибка.}
Последняя корректная часть рассуждения: если $p$ и $q$ пересекаются, то у них есть общая точка, и эта точка действительно принадлежит обеим плоскостям, поскольку обе плоскости содержат обе прямые.

Первый неверный переход возникает дальше: из того, что одна и та же точка принадлежит $\alpha$ и $\beta$, нельзя заключить, что $\alpha=\beta$. Две различные плоскости могут иметь общую точку, а если они пересекаются, то обычно имеют даже целую общую прямую.

\textbf{Почему это неверно.}
Для доказательства равенства двух плоскостей нужно установить, что обе они определяются одним и тем же набором геометрических объектов, который задаёт плоскость однозначно. В этой задаче таким набором являются не только точка пересечения, а две пересекающиеся прямые $p$ и $q$ целиком.

Именно две пересекающиеся прямые задают единственную плоскость. Это и есть ключевое свойство, которое связывает условия задачи с требуемым выводом.

\textcolor{blue!60!black}{\textbf{Ключевая идея:}} если две плоскости содержат одни и те же две пересекающиеся прямые, то обе обязаны быть той единственной плоскостью, которую эти прямые задают.

\textbf{Исправление от последнего правильного шага.}
Мы уже знаем, что прямые $p$ и $q$ пересекаются. Поэтому можно сразу применить теорему: через две пересекающиеся прямые проходит ровно одна плоскость. Поскольку обе плоскости $\alpha$ и $\beta$ содержат и $p$, и $q$, обе они являются этой единственной плоскостью.

\textbf{Полное правильное решение.}
Прямые $p$ и $q$ пересекаются. По теореме о задании плоскости двумя пересекающимися прямыми существует ровно одна плоскость, содержащая одновременно $p$ и $q$.

По условию $p\subset\alpha$ и $q\subset\alpha$, поэтому $\alpha$ является плоскостью, содержащей обе прямые $p$ и $q$.

Также по условию $p\subset\beta$ и $q\subset\beta$, поэтому $\beta$ тоже является плоскостью, содержащей обе прямые $p$ и $q$.

Но такая плоскость единственна. Следовательно,
\[
\alpha=\beta.
\]

\textbf{Проверка.}
Предположим, что $\alpha\ne\beta$. Тогда существовали бы две разные плоскости, каждая из которых содержит одновременно пересекающиеся прямые $p$ и $q$. Это противоречит теореме о единственности плоскости, проходящей через две пересекающиеся прямые. Значит, вывод $\alpha=\beta$ обязателен.

\textcolor{teal!60!black}{\textbf{Итог:}} в этой задаче нужно использовать не просто общую точку плоскостей, а тот факт, что обе плоскости содержат одни и те же две пересекающиеся прямые. Именно они однозначно задают плоскость.

\textbf{Задача на закрепление.}
Плоскости $\gamma$ и $\delta$ содержат две прямые $r$ и $s$, причём $r\cap s=\{T\}$. Докажите, что $\gamma=\delta$.

\newpage
\phantomsection\label{task:9}
\section*{Задание 9}

\textbf{Текст задания.}
В плоскости $\alpha$ лежат параллельные прямые $a$ и $b$. В плоскости $\beta$ лежит прямая $a$ и точка $M\in b$. Докажите, что $\alpha=\beta$.

\textbf{Ваше решение.}
Сначала отмечено, что параллельные прямые $a$ и $b$ задают единственную плоскость. Затем на прямой $a$ выбраны точки $A$ и $B$ и записано, что по аксиоме они задают единственную плоскость, после чего сделан вывод $\alpha=\beta$.

\textbf{Ваш ответ.}
$\alpha=\beta$.

\textcolor{red}{\textbf{Ошибка:}} две точки на одной прямой не задают единственную плоскость.

\textbf{Где именно возникает ошибка.}
Последний корректный шаг: параллельные прямые $a$ и $b$ действительно задают единственную плоскость; в условии этой плоскостью является $\alpha$.

Первый неверный шаг: после выбора двух точек $A$ и $B$ на прямой $a$ сказано, что эти точки задают единственную плоскость. Это неверно. Через одну прямую, а значит и через две точки $A$ и $B$ этой прямой, можно провести бесконечно много плоскостей.

\textbf{Почему это неверно.}
Чтобы получить единственную плоскость с помощью точек, нужны три точки, не лежащие на одной прямой. В этой задаче третья необходимая точка уже дана: это $M$.

Поскольку $a\parallel b$, прямые $a$ и $b$ не имеют общих точек. Следовательно, точка $M\in b$ не принадлежит $a$. Значит, если взять различные точки $A,B\in a$, то точки $A$, $B$ и $M$ не лежат на одной прямой. Именно они задают единственную плоскость.

\textcolor{blue!60!black}{\textbf{Ключевая идея:}} нужно использовать не только две точки прямой $a$, но и точку $M$, которая лежит вне $a$. Прямая и точка вне неё задают единственную плоскость.

\textbf{Исправление от последнего правильного шага.}
После того как установлено, что $a\parallel b$, замечаем: $M\in b$, поэтому $M\notin a$. Плоскость $\alpha$ содержит прямую $a$ и точку $M$, потому что $M\in b\subset\alpha$. Плоскость $\beta$ также содержит прямую $a$ и точку $M$ по условию. Но через прямую и точку, не лежащую на этой прямой, проходит ровно одна плоскость. Значит, $\alpha=\beta$.

\textbf{Полное правильное решение.}
По условию $a\parallel b$. Поэтому $a$ и $b$ не пересекаются.

Точка $M$ лежит на прямой $b$. Если бы $M$ лежала и на $a$, то прямые $a$ и $b$ имели бы общую точку $M$, что невозможно для параллельных прямых. Следовательно,
\[
M\notin a.
\]

Через прямую $a$ и точку $M$, не лежащую на этой прямой, проходит ровно одна плоскость.

Плоскость $\alpha$ содержит $a$ по условию. Кроме того, $M\in b$ и $b\subset\alpha$, поэтому $M\in\alpha$. Значит, $\alpha$ проходит через прямую $a$ и точку $M$.

Плоскость $\beta$ по условию также содержит прямую $a$ и точку $M$.

Обе плоскости проходят через одну и ту же прямую $a$ и одну и ту же точку $M\notin a$. Такая плоскость единственна, поэтому
\[
\alpha=\beta.
\]

\textbf{Проверка.}
Можно проверить вывод через три точки. Возьмём различные точки $A,B\in a$. Так как $M\notin a$, точки $A,B,M$ не лежат на одной прямой и задают единственную плоскость. И $\alpha$, и $\beta$ содержат все три точки, поэтому обе совпадают с этой единственной плоскостью.

\textcolor{teal!60!black}{\textbf{Итог:}} правильный вывод в Вашем решении получен, но в доказательстве нельзя заменять три неколлинеарные точки двумя точками одной прямой. Необходимую третью точку $M$ нужно явно использовать.

\textbf{Задача на закрепление.}
В плоскости $\gamma$ лежат параллельные прямые $c$ и $d$. Плоскость $\delta$ содержит прямую $c$ и точку $P\in d$. Докажите, что $\gamma=\delta$.

\newpage
\phantomsection\label{task:10}
\section*{Задание 10}

\textbf{Текст задания.}
Прямые $a$ и $b$ скрещиваются. На прямой $b$ выбраны различные точки $M$ и $N$. Пусть $\alpha$ --- плоскость, заданная прямой $a$ и точкой $M$, а $\beta$ --- плоскость, заданная прямой $a$ и точкой $N$. Докажите, что $\alpha\ne\beta$.

\textbf{Ваше решение.}
Записаны данные задачи и выполнен рисунок. Доказательство не приведено.

\textbf{Ваш ответ.}
Не указан.

\textcolor{red}{\textbf{Ошибка:}} отсутствует обязательное доказательство требуемого неравенства плоскостей.

\textbf{Где именно возникает ошибка.}
Данные записаны верно по смыслу: прямые $a$ и $b$ скрещиваются, а точки $M$ и $N$ лежат на $b$. На этом рассуждение останавливается. Первый недостающий шаг --- предположить совпадение плоскостей и вывести из этого противоречие со скрещиванием прямых.

\textbf{Почему этого недостаточно.}
Рисунок помогает представить конфигурацию, но не является доказательством. На рисунке плоскости могут выглядеть различными, однако требуется обосновать это только из условий задачи и известных свойств пространства.

Ключевое свойство скрещивающихся прямых: они не лежат в одной плоскости. Поэтому достаточно показать, что из предположения $\alpha=\beta$ следовало бы, что и $a$, и $b$ лежат в одной плоскости.

\textcolor{blue!60!black}{\textbf{Ключевая идея:}} если бы $\alpha$ и $\beta$ совпадали, общая плоскость содержала бы прямую $a$ и обе точки $M,N$. Но через две различные точки $M$ и $N$ проходит прямая $MN$, а поскольку обе точки лежат на $b$, имеем $MN=b$. Тогда общая плоскость содержала бы и $a$, и $b$, что невозможно для скрещивающихся прямых.

\textbf{Исправление.}
После записи данных нужно перейти к доказательству от противного. Предположим, что $\alpha=\beta$. Обозначим общую плоскость через $\pi$. Тогда $a\subset\pi$, $M\in\pi$ и $N\in\pi$. Так как $M\ne N$, вся прямая $MN$ лежит в $\pi$. Но $M,N\in b$, поэтому $MN=b$. Получаем $b\subset\pi$. Значит, $a$ и $b$ лежат в одной плоскости, что противоречит их скрещиванию.

\textbf{Полное правильное решение.}
Прямые $a$ и $b$ скрещиваются. Следовательно, они не лежат в одной плоскости.

Точки $M$ и $N$ различны и обе принадлежат прямой $b$.

Предположим противное: пусть
\[
\alpha=\beta.
\]
Обозначим эту общую плоскость через $\pi$.

По определению плоскости $\alpha$ она содержит прямую $a$ и точку $M$. Поэтому
\[
a\subset\pi,\qquad M\in\pi.
\]

По определению плоскости $\beta$ она содержит прямую $a$ и точку $N$. Поэтому
\[
N\in\pi.
\]

Итак, две различные точки $M$ и $N$ принадлежат плоскости $\pi$. Если две точки прямой принадлежат плоскости, то вся прямая, проходящая через эти точки, лежит в этой плоскости. Следовательно,
\[
MN\subset\pi.
\]

Но точки $M$ и $N$ лежат на прямой $b$ и различны, поэтому прямая $MN$ совпадает с $b$. Значит,
\[
b\subset\pi.
\]

Таким образом, одна плоскость $\pi$ содержит одновременно прямые $a$ и $b$. Это означает, что $a$ и $b$ лежат в одной плоскости, что противоречит условию о том, что они скрещиваются.

Следовательно, предположение о совпадении плоскостей неверно, и
\[
\alpha\ne\beta.
\]

\textbf{Проверка.}
Противоречие получено непосредственно с определением скрещивающихся прямых: если бы плоскости совпали, обе исходные прямые оказались бы компланарными. Значит, доказательство не использует свойства рисунка и зависит только от условий задачи.

\textcolor{teal!60!black}{\textbf{Итог:}} различие плоскостей доказывается от противного через прямую $MN=b$: две точки $M$ и $N$ в предполагаемой общей плоскости заставляют всю прямую $b$ лежать в ней, что невозможно вместе с $a$ при скрещивающихся прямых.

\textbf{Задача на закрепление.}
Прямые $c$ и $d$ скрещиваются. Через прямую $c$ и точку $P\in d$ проведена плоскость $\gamma$. Докажите, что никакая другая точка $Q\in d$, $Q\ne P$, не лежит в плоскости $\gamma$.

\vfill
\begin{center}
\small Лёвин Артём Александрович эксклюзивно для Ксении Клыковой
\end{center}

\end{document}
